% =============================================================================
% Appendix D --- Pulse Timing Details
% Grounded in Context 1 §4.6-4.7, Context 2 §5, Context 5 §1.2, and
% the actual weight_pulse_lut_table.txt in target/.
% =============================================================================

\chapter{Pulse Timing Details}
\label{app:pulse-timing}

This appendix collects the cycle-accurate specification of the pulse
machinery referenced throughout \Cref{ch:rtl}.

\section{Pulse Train Arithmetic}
\label{app:pulse-math}

The controller's pulse helpers parameterize every pulse phase by three
integers:

\begin{itemize}
    \item \(T\) --- active cycles per burst (package constants
          \sv{TREAD}, \sv{TPROG}, \sv{TERASE}, \sv{TINF}).
    \item \(N\) --- number of bursts per phase
          (\sv{PULSE\_NUM\_READ}, \sv{PULSE\_NUM\_PROG},
          \sv{PULSE\_NUM\_ERASE}, \sv{PULSE\_NUM\_INF}).
    \item \(G\) --- \signame{HIZ} gap cycles between adjacent bursts
          (\sv{PULSE\_GAP}).
\end{itemize}

\paragraph{Total length.} The package function
\sv{pulse\_train\_total(T, N, G)} returns the total cycle count
\(N \cdot T + (N-1) \cdot G\) (clipped to 8-bit unsigned).

\paragraph{Per-cycle activity.} \sv{pulse\_train\_active(cycle\_idx, T, N, G)}
returns 1 when \signame{cycle\_idx} falls inside an active burst and 0
when it falls inside a gap. The controller uses it to decide whether
\signame{pulses} takes the corresponding \sv{PULSE\_MODE\_*} or
\sv{PULSE\_MODE\_HIZ} on any given cycle.

\section{Per-State \signame{pulse\_total}}
\label{app:pulse-pertstate}

\Cref{tab:pulse-total-per-state} summarizes the value the combinational
block assigns to \signame{pulse\_total} depending on the main and
sub-state.

\begin{table}[htbp]
    \centering
    \caption{\signame{pulse\_total} per state.}
    \label{tab:pulse-total-per-state}
    \begin{tabularx}{\linewidth}{l l X}
        \toprule
        \textbf{Main state} & \textbf{Condition} & \textbf{\signame{pulse\_total}} \\
        \midrule
        \stateval{S\_READ} & --- &
            \sv{PULSE\_TOTAL\_READ} =
            \sv{pulse\_train\_total(TREAD, PULSE\_NUM\_READ, PULSE\_GAP)}. \\
        \stateval{S\_PROGRAM} & \sv{PROG\_WRITE}, \sv{USE\_WEIGHT\_PULSE\_LUT}=0 &
            \sv{pulse\_train\_total(TPROG, PULSE\_NUM\_PROG, PULSE\_GAP)}. \\
        \stateval{S\_PROGRAM} & \sv{PROG\_WRITE}, LUT on,
            \signame{prog\_retry\_cnt}~\(>\)~0 &
            Short retry train: \sv{Tp}~=~1, \sv{Np}~=~1 (one active cycle). \\
        \stateval{S\_PROGRAM} & \sv{PROG\_WRITE}, LUT on, first attempt &
            \sv{pulse\_lut\_macro\_repeat\_total(Rlut, Mmacro, PULSE\_GAP)},
            where \sv{Mmacro} is the non-LUT PROG total and
            \sv{Rlut}~=~\sv{weight\_pulse\_cycles\_lut[expected\_weight]}
            (minimum 1). \\
        \stateval{S\_ERASE} & \sv{ERASE\_PULSE} &
            \sv{PULSE\_TOTAL\_ERASE}. \\
        \stateval{S\_COMPUTE} & --- &
            \sv{PULSE\_TOTAL\_INF}, defined in the package as
            \(\max(8, \mathrm{TINF}\cdot N + \text{gaps})\). \\
        \bottomrule
    \end{tabularx}
\end{table}

\signame{pulse\_done} is
\sv{(pulse\_total > 0) \&\& (pulse\_cnt >= pulse\_total - 1)};
outside the active states above, \signame{pulse\_total}~=~0 and
\signame{pulse\_done}~=~0 always.

\section{LUT-Based Macro Repeat}
\label{app:pulse-lut-macro}

The function \sv{pulse\_lut\_macro\_repeat\_total(R, M, G)} concatenates
\(R\) copies of an inner macro of length \(M\) separated by
\(G\)-cycle gaps, giving a total of \(R\cdot M + (R-1)\cdot G\) cycles.
The companion \sv{pulse\_lut\_macro\_repeat\_active(cycle\_idx, R, M, G)}
returns 1 inside the macros and 0 inside the inter-macro gaps, so the
outgoing \signame{pulses} bus alternates between \sv{PULSE\_MODE\_PROG}
and \sv{PULSE\_MODE\_HIZ} \(R\) times.

\section{Weight Pulse LUT}
\label{app:pulse-lut-table}

\Cref{tab:weight-lut} lists the LUT contents shipped in
\sv{target/Controller/programming\_inputs/weight\_pulse\_lut.mem} (as
documented in \sv{weight\_pulse\_lut\_table.txt}). The row indexed by
the expected weight nibble gives the repeat count \sv{Rlut} used as the
first PROG shot for that cell.

\begin{table}[htbp]
    \centering
    \caption{Weight-indexed macro-repeat count \sv{Rlut} shipped in the
        project repository.}
    \label{tab:weight-lut}
    \begin{tabular}{c c c c c c c c c c c c c c c c c}
        \toprule
        Weight & 0 & 1 & 2 & 3 & 4 & 5 & 6 & 7 & 8 & 9 & 10 & 11 & 12 & 13 & 14 & 15 \\
        \sv{Rlut} & 2 & 5 & 3 & 6 & 7 & 4 & 8 & 5 & 9 & 6 & 10 & 8 & 4 & 11 & 6 & 12 \\
        \bottomrule
    \end{tabular}
\end{table}

Only the first PROG attempt reads this table; for every subsequent
retry, \signame{prog\_retry\_cnt}~\(>\)~0 forces the short single-cycle
train described above.

\section{Conceptual Pulse-Cycle Waveform}
\label{app:pulse-waveform}

\Cref{fig:app-pulse-cycle-waveform} shows a conceptual timing diagram
of one LUT-based PROG cycle, illustrating the relationship between
\signame{pulses}, \signame{pulse\_cnt}, \signame{pulse\_done}, and
\signame{prog\_state} during the \sv{PROG\_WRITE} phase.

\begin{figure}[htbp]
    \centering
    \includegraphics[width=0.95\linewidth]{fig_app_pulse_cycle_waveform}
    \caption[Conceptual PROG cycle waveform]{Conceptual LUT-based
        \sv{PROG\_WRITE} waveform. \signame{pulses} alternates between
        \sv{PULSE\_MODE\_PROG} and \sv{PULSE\_MODE\_HIZ} for \sv{Rlut}
        macro repetitions. \signame{pulse\_cnt} starts at 0 on entry
        to \sv{PROG\_WRITE} and increments every cycle;
        \signame{pulse\_done} rises on the last cycle of the last
        repetition, triggering the transition to \sv{PROG\_WAIT\_ACK}.}
    \label{fig:app-pulse-cycle-waveform}
\end{figure}
